\documentclass{article}
\usepackage[T1]{fontenc}
\usepackage{lmodern}
\usepackage{graphicx}
\usepackage{amsmath,amssymb}
\usepackage[a4paper,margin=2cm]{geometry}
\usepackage[absolute,overlay]{textpos}
\setlength{\TPHorizModule}{1mm}
\setlength{\TPVertModule}{1mm}
\usepackage[indonesian]{babel}
\usepackage{hyperref}
\setlength{\emergencystretch}{2em}
% unit: Halaman 31

\begin{document}

% === Halaman 31 (llm) ===

\textcolor{red}{Plainteks:}\vspace{2mm}
\begin{center}
\begin{tabular}{|p{0.6\textwidth}|}
\hline
Dinas Pendidikan Kota Ternate meminta kepada pihak sekolah dan orang tua siswa untuk jenjang pendidikan SD dan SMP se-Kota Ternate untuk melarang para siswa membawa permainan lato-lato yang sedang tren itu ke sekolah, karena akan mengganggu kegiatan belajar mengajar yang dinilai berbahaya sehingga mengantisipasi kecelakaan bagi anak di daerah itu.\hline
\end{tabular}
\end{center}
\vspace{10mm}
\textcolor{red}{Cipherteks:}\vspace{2mm}
\begin{center}
\begin{tabular}{|p{0.6\textwidth}|}
\hline
HAWFHZDOHAHANGOMKLGFCVWFLBOCPRKFGNHOFNIN \\ SPGNLHMPFVBEFWMVFWTBWRHSFZRWKFQMVHPAFWIK \\ DOHAHANGPFEWNFPNFKNLHMPFGLBAMGFCXKFQ MOCFV \\ AVKANIAVHSFZRNCODGAFOHYUVGWFOGFXEGFMLFWIT \\ HEFWTBVCPAYQOBLHMPFVBPVADOVWGNGFWOCKARVKA \\ RQOBRGGFFWCHFVGVYUDORVGVYCFWABAPVSVGCHGCI \\ DRFIFH BAPVQRVOCKAFWSGHSNIHSOBDHFVNGFWDGRGF \\ WGNHANFCVIDGSWZ \hline
\end{tabular}
\end{center}

\end{document}
